\chapter{Lifting the actual axioms, proof rules and targets}
\label{tx:chap:assemble-transfer}
All the finite-algebra obligations have now been supplied. We return
to the original question $\Gamma_n\vdash G_{b,n}$: can every requested
target be derived from the given axioms within degree $d$? We verify each part of the contract
in Chapter~\ref{tx:chap:contract}. The point of doing the assembly slowly
is to make clear that a prescribed target, rather than an arbitrary
lifted vector, is transported.

\section{The theorem we can now state without hidden assumptions}
Let $\Gamma_n$ be a fixed uniform axiom template and let $G_{b,n}$ be
a fixed uniform target family. Put
\[
 R(d)=2kd+s+t,\qquad n_*=\max\{n_0,1,kD_G+t\}.
\]
Here $k,s,t$ bound the label supports of a variable, an axiom index and
a target index; $D_G$ bounds target degree. The rank $n_0$ is the axiom
stability threshold. These are fixed data; only $d$ is a proposed proof degree.

\begin{theorem}[Actual target transfer]\label{tx:thm:transfer}
Over any field of fixed prime characteristic $p$, suppose
\[
 q=p^a>R(d),\qquad n\ge\max\{n_*,2R(d)-1\},\qquad N=n+q.
\]
For either NS or ordinary PC, every prescribed rank-$N$ target has a
derivation of degree at most $d$ if and only if every prescribed rank-$n$
target does. The assertion holds for every $d\ge0$, with the separate
zero-degree interpretation given below. It also holds for any
representative of one fixed global target orbit at each eligible rank.
For a fully invariant target use the sharper budget $kd+s$ and start
$\max\{n_0,1,kD_G\}$.
\end{theorem}

The theorem concerns original-field polynomials and actual derivations.
It does not require invariant certificates, does not assume that the
axioms are inconsistent, and does not add the target to them. If a
stabilised target degree exceeds $d$, the threshold predicate is simply
false. The empty target request is true at every degree.

\section{Transporting the entire consequence space}
Fix a positive degree budget $d$ and write $W_N=P_{N,\le d}$
for the ordinary polynomial space of degree at most $d$. Recall
\[
 I_d(N)=\Span_K\{\mu f_{a,N}:\deg\mu+\deg f_{a,N}\le d\},
 \qquad C_d(N)=V_\infty(N).
\]
The first space records NS identities of total-product degree at most
$d$, with $\mu$ ranging over monomials; the second is the closure of the usable axioms under the legal
ordinary-PC rules (Proposition~\ref{tx:prop:closure}). Boolean equations,
when present, remain inputs in this polynomial space.
We first transport the spaces $I_d(N)$ and $C_d(N)$, without choosing
any target. This separates the proof-system calculation from the
question of whether a particular uniform expression belongs to it.
The operation $\Del(M)=\ker z_M/(\ker z_M\cap\im z_M)$ deletes the
block through the label action; the common projector represents its
classes, and $\Lift=e\iota$ lifts smaller-rank vectors. Exactness means
that kernels, images and quotients are preserved by this operation.

Concretely, a deletion class is represented by a vector $v$ with
$z_Mv=0$. Two such vectors represent the same class if their difference
lies in $\ker z_M\cap\im z_M$. The projector provides representatives
which can be compared with smaller-rank vectors. We use the actual
inclusions and maps between these classes, not just equal dimensions
of the resulting spaces (Chapter~\ref{tx:chap:deletion}).

For NS the certificate map has source basis the pairs $(\mu,a)$ with
\[
 \deg\mu+\deg f_{a,N}\le d,
 \qquad (\mu,a)\longmapsto\mu f_{a,N}\in W_N.
\]
Its source has support at most $kd+s$, and $W_N$ has support at most
$kd$. The matrix entry at a target monomial $\nu$ is zero unless
$\mu\mid\nu$, and otherwise is $c(a,\nu/\mu)$. These are the actual
rank-independent invariant entries of the certificate map. The
stability threshold $n_0$ makes each axiom degree independent of rank.
Exactness and matrix naturality therefore identify the embedded image
$\Del(I_d(N))$ with $I_d(n)$.

For PC let $X_N=K[\mathcal V_N]$ and $L_N=P_{N,\le d-1}$.
The tensor product records a variable and a lower-degree line together:
a pure tensor $x_v\otimes f$ is sent to their polynomial product,
and sums of pure tensors record sums of such products. The degree
restriction is imposed on $f$ before multiplication. This is the
variable-labelled register of Chapter~\ref{tx:chap:contract}.
The multiplication map
\[
       m_N:X_N\otimes L_N\longrightarrow W_N,\qquad
                 x_v\otimes f\longmapsto x_vf
\]
has stable entries. Its source is supported on at most
$k+k(d-1)=kd$ labels. The usable axioms give the image $V_0(N)$
of their coefficient map. The recurrence is
\[
 V_{j+1}(N)=V_j(N)+\im\bigl(m_N|_{X_N\otimes(V_j(N)\cap L_N)}\bigr).
\]
Every stage is a submodule of $W_N$, and every multiplication domain
is a submodule of the fixed support-bounded space $X_N\otimes L_N$.
Suppose stage $j$ has been identified at both ranks. Put
$U_N=V_j(N)\cap L_N$ and $U_n=V_j(n)\cap L_n$.
The intersection is the kernel of $(v,l)\mapsto v-l$ on
$V_j(N)\oplus L_N$, so exactness identifies $\Del(U_N)$ with
$U_n$, including their inclusions $i_N:U_N\hookrightarrow L_N$
and $i_n:U_n\hookrightarrow L_n$.

Write $m'_N=m_N\circ(1\otimes i_N)$ and
$m'_n=m_n\circ(1\otimes i_n)$ for the restricted multiplication
maps. The required compatibility is the commutative square
\begin{equation}\label{tx:eq:restricted-multiplication}
\begin{CD}
 \Del(X_N\otimes U_N) @>{\Del(m'_N)}>> \Del(W_N)\\
 @V{\simeq}VV @V{\simeq}VV\\
 X_n\otimes U_n @>{m'_n}>> W_n.
\end{CD}
\end{equation}
The left identification is \eqref{tx:eq:tensordelete}, natural in
the second factor: it sends $\Del(1\otimes i_N)$ to
$1\otimes i_n$. The ambient multiplication map $m_N$ is sent to
$m_n$ by stable matrix naturality. Composing these two statements
proves the square for the restriction. Exactness now transports its
image and the sum with $V_j$, giving stage $j+1$. The initial stage
is the actual usable-axiom image, so induction identifies the embedded
$\Del(V_j(N))$ with $V_j(n)$ for every $j$.

Both ascending chains stabilise because their ambient spaces are finite
dimensional. Choose a stage beyond both stabilisation points. We obtain
\begin{equation}\label{tx:eq:consequence-transport}
 \Del(I_d(N))=I_d(n),\qquad \Del(C_d(N))=C_d(n),
\end{equation}
where equality means equality after the natural identification
$\Del(W_N)\simeq W_n$. This transports every cancellation in a PC
proof, including a low-degree result subsequently multiplied by a
variable. No bound on proof length and no symmetry of individual proof
lines was used.

\section{Why deletion also reflects target membership}
Preserving an image is not by itself enough: an exact operation can
annihilate a nonzero quotient. The next lemma supplies the reflection
needed for a general target source.

The notation $S'/S$ means vectors in $S'$ considered equal when their
difference lies in $S$; it is a \emph{subquotient} of the ambient space.
In the lemma, $\mathcal C_{2u}(N)$ is the common-annihilator class of
Chapter~\ref{tx:chap:deletion}, not a PC consequence space.
Here $u$ is an integer support bound, unrelated to the top-power
operator $u=z^{p-1}$ used in the witness construction.
The doubled budget has a concrete purpose: a basis pair in
$P\otimes P$ can use the union of two support-$u$ label sets.
The annihilator class is not being assumed closed under arbitrary
tensor products.

\begin{lemma}[Nonzero subquotients survive]\label{tx:lem:detect}
Let $P=K[\mathcal U_N]$ have support at most $u$. Suppose deletion is
exact on $\mathcal C_{2u}(N)$, its block size exceeds $2u$, and the
natural identifications \eqref{tx:eq:permdelete}--\eqref{tx:eq:tensordelete}
hold. If $S\subsetneq S'\subseteq P$ are invariant subspaces, then
$\Del(S'/S)\ne0$.
\end{lemma}
\begin{proof}
The permutation basis gives a nondegenerate dot product on $P$,
equivariant under $S_N$. It induces a pairing
\[
                         S^\perp\otimes(S'/S)\longrightarrow K.
\]
This pairing is surjective: for $v\in S'\setminus S$, a linear
functional vanishing on $S$ can take value one at $v$; the dot product
represents that functional by an element of $S^\perp$.
Let $Y=S'/S$. Tensoring over a field is exact, so the displayed
source lies in $P\otimes Y$, itself a subquotient of $P\otimes P$.
The last permutation module has support at most $2u$, justifying
every use of exactness. If $\Del(Y)=0$, permutation-first-factor tensor
compatibility gives $\Del(P\otimes Y)=0$. Deleting the inclusion
$S^\perp\otimes Y\hookrightarrow P\otimes Y$ gives an injection
\[
                  \Del(S^\perp\otimes Y)\hookrightarrow0,
\]
so the pairing source also deletes to zero. Exactness applied to the
surjective pairing then gives $0\twoheadrightarrow\Del(K)\simeq K$,
which is impossible because $K\ne0$.

We use this surjective pairing, not a trace or an averaged inner
product. In characteristic $p$, composing a pairing with coevaluation
can yield a vanishing dimension scalar and would not prove the lemma.
\end{proof}

In the detector proof, $S^\perp$ consists of vectors whose dot product
with every vector of $S$ is zero. Thus pairing $h\in S^\perp$ with
$v+S\in S'/S$ is well-defined: replacing $v$ by $v+s$ changes the
value by $\langle h,s\rangle=0$. The tensor product is the linear
source of this bilinear pairing. Surjectivity says that some such
pair has value one, rather than that a dimension scalar is nonzero.

\begin{lemma}[An invariant target needs no doubled budget]\label{tx:lem:constant}
Let $j:K\to W$ and $H\subseteq W$ be equivariant and belong to the
exactness class. Then $j(1)\in H$ if and only if its deleted image
belongs to $\Del(H)\subseteq\Del(W)$.
\end{lemma}
\begin{proof}
Membership is the vanishing of $K\xrightarrow{j}W\to W/H$.
Its image is either zero or isomorphic to the simple trivial module
$K$. Exactness preserves this image, and $\Del(K)\simeq K\ne0$.
Thus the composite vanishes exactly when its deleted composite does.
The vector $j(1)$ may be any invariant target, not just the constant
polynomial one.
\end{proof}

\section{The actual uniform target map}\label{tx:sec:target-transfer}
For indexed targets, let $\mathcal B$ be the fixed index basis and
let $\gamma(b,\mu)$ be its stable coefficient rule. Write $\delta$ for the largest stabilised target degree, taking
$\delta=0$ for an empty request. At a rank where
all target degrees are at most $d$, define
\[
 J_N:K[\mathcal B_N]\longrightarrow W_N,\qquad
                [b]\longmapsto G_{b,N}.
\]
Here $J_N$ is a linear target map. It is distinct from the group-algebra
annihilator ideal $J_r(N)$ used to construct the common projector.
Its entry in row $\mu$, column $b$ is precisely $\gamma(b,\mu)$.
Equivariance and rank independence are the hypotheses of the natural
matrix identification in Theorem~\ref{tx:thm:deletion}. Consequently
$\Del(J_N)$ is the actual smaller target map $J_n$.
This statement concerns the prescribed full polynomials $G_{b,n}$;
we do not replace a target by its degree-$d$ truncation.

Let $H_N$ be either $I_d(N)$ or $C_d(N)$. Every target belongs to
$H_N$ exactly when $\im J_N\subseteq H_N$. Put
\[
          S_N=H_N+\im J_N,\qquad Y_N=S_N/H_N.
\]
The condition is $Y_N=0$. Exactness, target-map naturality and
\eqref{tx:eq:consequence-transport} identify its deletion with
\[
          Y_n=(H_n+\im J_n)/H_n.
\]
The polynomial permutation basis of $W_N$ has support at most $kd$.
Lemma~\ref{tx:lem:detect}, applied with the original monomial support bound,
reflects vanishing of $Y_N$ once the budget covers $2kd$.
The affine budget
\[
                            R(d)=2kd+s+t
\]
also covers the certificate source ($kd+s$), the target source ($t$)
and every multiplication domain. It therefore proves membership in
both directions. The target source need not be simple; reflection
comes from a subquotient of the ambient polynomial target space.
Doubling $R(d)$ once again would be unnecessary: the pairing uses
$W_N\otimes W_N$, whose support is $2kd$.

\begin{proof}[Proof of Theorem~\ref{tx:thm:transfer}]
First extend scalars to an algebraic closure, as permitted by
Lemma~\ref{tx:lem:field}. The degree of every indexed target stabilises
beyond $n_*=\max\{n_0,1,kD_G+t\}$. If $d$ is smaller than the
largest stable target degree, the all-target predicate is false at both
ranks. Otherwise the full map $J$ just described applies.
Take $N=n+q$, $q=p^a>R(d)$ and retained rank
$n\ge\max\{n_*,2R(d)-1\}$. Exact deletion transports $H_N$, $J_N$
and their sum and quotient. The detection lemma yields
\[
 \bigl(\forall b\in\mathcal B_N,\ G_{b,N}\in H_N\bigr)
 \quad\Longleftrightarrow\quad
 \bigl(\forall b\in\mathcal B_n,\ G_{b,n}\in H_n\bigr).
\]
Field descent holds for each actual target separately, so the equivalence
holds over the original field. Repeated steps connect any eligible
ranks in a common residue class and prove the stated periodicity.
For one fixed global target orbit, replace $\mathcal B$ by that orbit;
at eligible ranks, equivariance makes every representative equivalent.

At degree zero the stabilised targets are either all zero, or are
scalars with at least one nonzero value. In the latter case every target
is derivable in degree zero exactly when the original axioms include a
nonzero constant. If a target has positive stable degree, degree-zero
all-target membership fails. These are fixed predicates beyond $n_*$
and give the separate zero-level assertion.

For a fully invariant target use $j_N:K\to W_N$, $j_N(1)=G_N$,
in place of $J_N$. Lemma~\ref{tx:lem:constant} reflects membership without
the tensor pairing. The budget reduces to $kd+s$, with stable onset
$\max\{n_0,1,kD_G\}$. Taking $G_N=1$ gives the sharper refutation
special case. No target has been added to the axiom system.
\end{proof}


\section{From one deletion step to a residue class}
Write $h_S(n)$ for the maximum least degree of the requested targets
in system $S$. Recall the safe onset $N(d)$ and chosen $p$-power period
$Q(d)$ from \eqref{tx:eq:NQ}; their separate zero-degree values are
$N(0)=n_*$ and $Q(0)=1$.
Fix $d$ and an eligible power $q$. If two sizes at or above the chosen
onset have the same residue modulo $q$, repeatedly subtract $q$ from
the larger until reaching the smaller. Every retained size in those
steps remains eligible. Applying Theorem~\ref{tx:thm:transfer} in both
directions proves threshold periodicity.

To preserve an exact minimum degree rather than just the assertion
$h(n)\le d$, repeat the argument at every lower threshold. The periods
are powers of the same prime and divide the chosen degree-$d$ period;
zero degree has a stable separate predicate. The vector of all those
threshold answers determines the truncated value
$\min\{h(n),d+1\}$. This completes the input used in
Theorem~\ref{tx:thm:profiles} for the four profiles.

\section{The complete lower-bound argument in a paragraph}
Fix a minimum finite degree $d$ occurring at a size beyond its own onset.
The two-way threshold theorem makes that exact degree recur along a
prime-power progression. Therefore a degree which is not allowed to recur
must occur before its onset. Since the onset is $O(d)$, those isolated
occurrences must have degree $\Omega(n)$. On T3 hard residues the same
contradiction excludes every recurring easy level. On T4 we instead use
unbounded recurring levels and the short period to bring a representative
down to size $O(d)$, producing finite linear peaks. No iteration ever
reduces the absolute degree: it is the availability of recurrence and
its onset which produce the gap.
